\section{Related Work}

In this section,
we review prior deterministic threshold signature schemes in the literature.
We review randomized threshold Schnorr signature schemes in
Section~\ref{frost-related-work}.

The BLS signature scheme~\cite{BonehLS04} is deterministic,
because the signature is of the form $z \gets H(m)^\sk$.
Likewise, threshold BLS signatures are similarly deterministic~\cite{Boldyreva03}.
However, such schemes are often not viable in a practical setting that requires backwards compatibility with Schnorr signature verification,
or due to the requirement of pairings.

Deterministic threshold Schnorr signatures have been described in the literature,
but rely on heavyweight zero-knowledge proofs to demonstrate that each
participant followed the protocol correctly.
The challenge of deterministic threshold Schnorr signatures is that of \emph{verifiability},
because if even one participant deviates from the protocol and chooses its nonce randomly,
a complete key-recovery attack is possible.

Nick et al.~\cite{NickRSW20} define a threshold Schnorr $n$-of-$n$ multisignature
that uses SNARKs to demonstrate in zero-knowledge that
a participant generated its nonce using a PRF correctly with respect to a pre-established keypair.
However, the authors report the computational overhead of at least 943\,ms for a single execution,
% \ian{The 943\,ms figure is just for ZKP creation; there's more for verification, and all the non-ZKP bits.}
independent of the number of signers,
due to the overhead of proving the PRF was evaluated correctly in zero-knowledge.
In particular, to instantiate the PRF,
a non-algebraic cryptographic hash function $H : \zo \rightarrow \GG$ must be used,
along with a regular function $f : \GG \rightarrow \Zq$.
The complete PRF $F_\sk$ is defined by $F_\sk(x) = f( \sk \cdot H(x))$,
where $\sk$ is a PRF key.
Then, signers generate a Bulletproof~\cite{BunzBBPWM18} to prove in zero-knowledge that the nonce was derived with respect to $F$, $\sk$, and some input $H(x)$.

Bronte et al.~\cite{BonteST21} define an MPC protocol to partition the functionality of EdDSA,
which itself is deterministic.
However, their MPC protocol is randomized and requires multiple rounds of interaction,
and therefore is stateful.

Garillot et al.~\cite{GarillotKMN21} similarly rely on parties sampling and committing to a PRF key at the time of key generation,
but instead make use of the Zero-Knowledge from Garbled Circuits (ZKGC) paradigm~\cite{JawurekKO13} to demonstrate that the PRF was derived correctly.
The specific function that is garbled is $C(x) = \phi(F(x))$,
where $F$ is a boolean circuit such as AES,
and $\phi$ is standard exponentiation (or curve multiplication).
The authors give efficiency estimates for a 256-bit curve  and using SHA-512 as
the PRF $F$;
the performance overhead is dominated by performing $132,000$ AES invocations and $256$ additions in $\Zq$ per proof generated and verified.
For a signing invocation involving $t$ parties,
each party must then perform $256t$ field operations and $132,000t$ AES invocations,
accounting for each proof a signer generates and the $t-1$ proof verifications for all other signers.

Kondi, Orlandi, and Roy~\cite{KondiOR23} take a different approach,
and define a two-round stateless and deterministic two-party Schnorr signature scheme using pseudorandom correlation functions (PCFs)~\cite{BoyleCGIKS20} as a building block.
In particular,
their scheme employs a Paillier-based PRF~\cite{OrlandiSY21},
but additionally define a verification mechanism to ensure that parties honestly followed the protocol.
However,
their scheme is restricted to the two-party setting,
and as written, cannot be extended to the general $(t,n)$ threshold setting.
In particular, the PCF used as a building block by their scheme assumes only two parties.
Extending their scheme to any $(t,n)$ setting requires designing a new $n$-party PCF~\cite{KPC24}.

\paragraph{Honest Majority Threshold Signatures.}
Honest majority threshold schemes have been proposed in the literature as a means to achieve properties that are either impossible,
or require higher performance overhead,
in the honest minority setting.
Notably, honest majority schemes have been demonstrated to achieve robustness or to circumvent requiring the use of heavyweight zero-knowledge proofs~\cite{CanettiGGMP20}.

Gennaro et al.~\cite{GennaroJKR96} use error correcting codes in the honest majority setting to achieve a robust threshold DSS scheme.
Similarly, Ruffing et al.~\cite{RuffingRJSS22} present a wrapper to the FROST threshold signature scheme to achieve robustness,
in the honest majority setting.

In the randomized threshold ECDSA setting,
Damg{\aa}rd et al.~\cite{DamgardJNPO22},
Doerner et al.~\cite{DoernerKLS23},
and Delskov~\cite{DalskovOKSS20}
show how  the pseudorandom secret sharing scheme by Cramer et al.~\cite{CramerDI05} can be employed as a sub-protocol for improved efficiency.
However, our work targets the setting of deterministic threshold Schnorr signatures.

\paragraph{Distributed VRFs.}
Galindo et al.~\cite{GalindoLOW21} give a formalization of distributed VRFs and their security notions.
While our notions for a Verifiable Pseudorandom Secret Sharing (VPSS) scheme are similar,
our definition for a VPSS does not require that players' outputs are accompanied by a zero-knowledge proof that the protocol was performed correctly.
Instead,
the VPSS verification function \emph{collectively} verifies all parties' outputs.


\paragraph{Pseudorandom Secret Sharing.}
Our VPSS construction \DVRFOne builds on the pseudorandom secret sharing scheme by Cramer et al.~\cite{CramerDI05}.
Note that while Cramer et al.\ additionally define a Non-Interactive\ Verifiable Secret-Sharing (NIVSS)  variant,
their construction assumes that players output shares in the clear,
and requires an honest supermajority so that the secret can be recovered.
In our case, players verify the correctness of their shares with respect to other players' public commitments,
and requires only honest majority (as players simply need to determine if any other party deviated from the protocol).

While we are the first to do so in a threshold Schnorr signature setting,
pseudorandom secret sharing has been used as a building block in other other threshold settings.
For example,
Jarecki, Krawczyk, and Resch~\cite{JareckiKR18} define a threshold Oblivious PRF which likewise builds upon pseudorandom secret sharing.

